%==============================================================================
% Paper F -- skeleton (bullets, to be polished into prose)
% Working title: "B+3.1: A Streaming, Delete-as-You-Go Whole-VM
%   Snapshot-Capture System"
% Role: this is the SYSTEMS / architecture paper for the capture engine. It
%   deliberately COMPLEMENTS Paper A (the signal paper): Paper A defines the
%   Active Page Fraction (APF) signal and its cost model; Paper F goes deep on
%   the capture machinery that produces it -- the asynchronous off-critical-path
%   helper, the delete-as-you-go transient-disk bound, the cell-end completeness
%   barrier, and the REJECTED concurrent multi-VM path.
% Target venue: systems / virtualization (e.g., a VEE/middleware-style venue or
%   a systems workshop).
%
% Build: pdflatex skeleton.tex   (needs IEEEtran; if unavailable, change the
%   \documentclass line below to \documentclass[11pt]{article} and delete the
%   \IEEEauthorblock* macros.)
%==============================================================================
\documentclass[conference]{IEEEtran}
\usepackage{amsmath,booktabs,siunitx,url}

\title{B+3.1: A Streaming, Delete-as-You-Go\\
Whole-VM Snapshot-Capture System}

\author{\IEEEauthorblockN{TODO Author list}
\IEEEauthorblockA{TODO Affiliation\\
\texttt{TODO@email}}}

\begin{document}
\maketitle

%------------------------------------------------------------------
\begin{abstract}
% TODO: polish into 150-200 words of prose. Bullet stub:
\begin{itemize}
  \item Problem: repeatedly snapshotting whole guest memory to compute a
        page-change signal is storage- and latency-bound. A naive
        ``snapshot-then-analyze-then-keep'' loop retains one \SI{1}{\gibi\byte}
        dump per snapshot and blocks the producer while the diff runs.
  \item System: \emph{B+3.1}, a capture engine in which the on-critical-path
        work is only the hypervisor snapshot primitive (\texttt{virsh suspend}
        $\rightarrow$ QEMU \texttt{pmemsave} $\rightarrow$ \texttt{virsh
        resume}), while an \emph{asynchronous} helper computes the page-change
        signal off the critical path, appends one trajectory record, and
        deletes the prior dump (\emph{delete-as-you-go}).
  \item Guarantees: a per-pair acknowledgement file plus a \emph{cell-end
        barrier} guarantee trajectory completeness before any downstream metric
        is computed; structured helper exit codes surface silent failures.
  \item Design target: delete-as-you-go is intended to bound transient disk to
        roughly \SI{2}{}--\SI{3}{\gibi\byte} regardless of trajectory length.
        % HONESTY: the v3 campaign ran with keep_dumps enabled, which DISABLES
        % the delete step. The 2-3 GiB cap is therefore a DESIGN TARGET, not a
        % measured result. See Limitations and Sec. evaluation.
  \item Negative result: concurrent multi-VM \texttt{pmemsave} was evaluated and
        \emph{rejected} -- a single shared SSD serializes the dump writes, so
        concurrency does not raise throughput and worsens \texttt{pmemsave}
        drift while breaking cell independence.
  \item Measured cost: \texttt{pmemsave} dominates the snapshot cycle (median
        \SI{7.06}{\second}, a median \SI{81}{\percent} of the cycle); suspend
        ($\approx$\SI{0.35}{\second}) and resume ($\approx$\SI{0.05}{\second})
        are negligible.
\end{itemize}
\end{abstract}

%------------------------------------------------------------------
\section{Introduction}
% Proposed introduction (bullets):
\begin{itemize}
  \item Hypervisor-external whole-VM memory capture writes a flat image of all
        guest physical memory per sample, then derives a change signal between
        consecutive images. The signal itself is the subject of Paper A; this
        paper is about the \emph{capture engine} that makes such repeated
        capture feasible on a single commodity host.
  \item Two pressures dominate the engine design:
    \begin{itemize}
      \item \textbf{Disk.} Each dump is \SI{1}{\gibi\byte}
            (guest RAM \SI{1024}{\mebi\byte}; page size \SI{4096}{\byte};
            $N = 262{,}144$ pages). Retaining every dump for even one long cell
            costs tens of \si{\gibi\byte}; retaining a whole campaign is
            multi-terabyte and infeasible on the available host
            (see Paper A's storage argument and the keep-dumps audit).
      \item \textbf{Latency.} A synchronous ``diff after each snapshot'' design
            (the B+3 baseline) blocks the producer for $\approx$\SI{1}{\second}
            per snapshot while it reads two \SI{1}{\gibi\byte} dumps; at small
            requested intervals this multiplies the cycle and inflates the
            guest pause fraction.
    \end{itemize}
  \item \textbf{This paper} presents B+3.1, which resolves both: the
        page-change helper runs \emph{asynchronously} (off the critical path),
        and a \emph{delete-as-you-go} mode is designed to hold only a small
        number of dumps on disk at once. A \emph{cell-end barrier} (waits for
        all helper acks) guarantees the trajectory is complete before metrics.
  \item \textbf{Honesty up front.} The v3 production campaign ran with
        \texttt{keep\_dumps} enabled, so the delete step was disabled and the
        \SI{2}{}--\SI{3}{\gibi\byte} transient cap was \emph{not} exercised;
        observed transient disk instead scaled with trajectory length. Closing
        this gap is the paper's biggest open item (R-3), deferred to the
        Plan~05 snapshot-throughput proposal.
        % HONESTY (prominent): keep_dumps=true on 125/128 ok cells; delete-mode
        % unexercised. The cap is a design target. Do not present it as measured.
  \item \textbf{Contributions} (expanded in Sec.~\ref{sec:contrib}):
    \begin{itemize}
      \item C1: a streaming capture architecture whose on-critical-path cost is
            only the hypervisor snapshot, with the change-signal helper moved
            off-path (Δ-1, async helper + cell-end barrier).
      \item C2: a delete-as-you-go transient-disk design (a \emph{target}
            bound, not yet validated) plus a per-pair ack / structured
            exit-code resilience layer (Δ-2) and a seq-numbered, sentinel-closed
            trajectory with a completeness validator (Δ-3, Δ-4).
      \item C3: a measured snapshot-cost characterization establishing
            \texttt{pmemsave} as the binding cost, and a documented
            \emph{negative result} rejecting concurrent multi-VM
            \texttt{pmemsave}.
    \end{itemize}
\end{itemize}

%------------------------------------------------------------------
\section{Background and Related Work}
\begin{itemize}
  \item Hypervisor snapshot / checkpoint primitives: the capture uses the QEMU
        monitor \texttt{pmemsave} together with libvirt \texttt{virsh
        suspend}/\texttt{resume}~\cite{TODO_libvirt_pmemsave}. TODO: cite
        snapshot/checkpoint cost studies to frame the \texttt{pmemsave}
        overhead.
  \item VM introspection / memory forensics~\cite{garfinkel2003vmi}: external
        observation of guest memory; this paper supplies the capture substrate
        such introspection runs on. TODO: position against live-migration
        dirty-page tracking.
  \item Working-set estimation / dirty-bit sampling~\cite{TODO_working_set_estimation}:
        hypervisors already track dirtied pages for migration; the present
        engine is content-external (it diffs full dumps) rather than
        dirty-bit-driven. Relevant as the natural future optimization that
        would remove the guest pause.
  \item Memory deduplication / same-page merging~\cite{TODO_ksm_dedup}: shares
        the page-granular content view; cited only for contrast with the
        capture model here.
  \item Asynchronous producer/consumer and write-back / page-cache effects:
        the engine's timing is sensitive to dirty-page writeback pressure.
        % TODO: add citation for OS page-cache / writeback behavior. The
        % flush <-> self-clean coupling (Paper E) is the in-project evidence.
        Forward reference: Paper E (timing instrumentation) characterizes the
        page-cache coupling and the interval-inertness of the snapshot cycle.
  \item TODO: cite a representative prior streaming-/online-diff system if one
        exists, to position delete-as-you-go.
\end{itemize}

%------------------------------------------------------------------
\section{Architecture}
\label{sec:arch}
\begin{itemize}
  \item \textbf{Capture primitive (on the critical path).} Per snapshot cycle
        the producer issues \texttt{virsh suspend}, waits for
        \texttt{domstate==paused}, invokes the QEMU monitor
        \texttt{pmemsave(val:0, size:ramBytes, file)} to write a flat RAW image
        (\texttt{memory\_dump-<ts>.raw}, exactly \SI{1}{\gibi\byte}), then
        \texttt{virsh resume}~\cite{TODO_libvirt_pmemsave}. This is the only
        guest-visible (pausing) work in B+3.1.
  \item \textbf{Asynchronous page-change helper (off the critical path).} After
        \texttt{pmemsave} returns, the producer spawns
        \texttt{plan02\_apf\_helper.py \&} in the background (the trailing
        \texttt{\&} backgrounds it) and proceeds immediately to
        \texttt{virsh resume} and the next sleep. The helper:
    \begin{itemize}
      \item maps both dumps as \texttt{uint8} via \texttt{numpy.memmap},
            reshapes to $[N,4096]$, and computes the page-change signal as a
            byte-wise per-page \emph{differ} reduction (XOR-equivalent
            \texttt{(a != b).any(axis=1)}), $\approx$\SI{1}{\second} wall-clock
            on the host SSD;
      \item appends exactly one record to the shared
            \texttt{apf\_trajectory.jsonl} via an atomic \texttt{O\_APPEND}
            write (line $<$ \texttt{PIPE\_BUF}, so concurrent appends do not
            interleave);
      \item deletes the \emph{previous} \texttt{.raw} (delete-as-you-go);
      \item writes a per-pair \texttt{.apf\_done} acknowledgement file via
            tempfile + atomic \texttt{rename}.
    \end{itemize}
    % HONESTY: the delete step is the line that keep_dumps disables. With
    % keep_dumps=true the helper still computes the signal and writes the
    % record + ack, but DOES NOT delete prev; dumps are retained until a single
    % end-of-cell sweep.
  \item \textbf{Why async (Δ-1).} In the synchronous baseline (B+3) the producer
        blocks $\approx$\SI{1}{\second}/snapshot on the helper. Async overlaps
        the helper's compute with the next cycle's snapshot work, dropping the
        effective per-cycle cost from $\approx$\SI{2.5}{\second} (sync) to
        $\approx$\SI{1.5}{\second} (async) as reported in the B+3.1 proposal.
        % HONESTY: 2.5 s -> 1.5 s is the proposal's reported on-path reduction
        % for this design; cross-check against Paper E's measured cycle numbers
        % before publication (the proposal's ~1 s helper and ~1.5 s host_dt are
        % from the timing-instrumentation regime, not the production v3 run).
  \item \textbf{Transient-disk model (delete-as-you-go).} Steady state holds
        \texttt{prev}~+~\texttt{curr} ($\approx$\SI{2}{\gibi\byte}) plus at most
        one in-flight helper's dump, i.e. a target peak of
        $\approx$\SI{3}{\gibi\byte} \emph{independent of trajectory length}.
        % HONESTY (R-3): this bound is a DESIGN TARGET. It was NOT exercised in
        % v3 (keep_dumps on). At the host level the v3 run showed a 59.5 GiB
        % free-disk swing because deletion was deferred to block-level sweeps.
        % Validation of the cap belongs to the Plan 05 snapshot-throughput
        % proposal.
  \item \textbf{Resilience: ack files + structured exit codes (Δ-2).} The
        helper returns one of five exit codes -- 0 ok, 1 transient io, 2 numpy
        missing, 3 disk full, 4 corrupt/zero-byte dump -- and the ack file is
        the last action before exit, so its presence is proof the pair
        completed. A missing ack is recorded as a gap. (Producer reactions:
        continue / log+continue / abort-cell after repeated code-2 / backpressure
        / mark-gap.)
  \item \textbf{Completeness: seq-numbered JSONL + final sentinel (Δ-3).} Each
        record carries an integer \texttt{seq}; under async, helpers can finish
        out of order, so the reader sorts by \texttt{seq} and checks for a
        contiguous range. The producer appends a final sentinel
        (\texttt{\{"final":true, "n\_pairs\_expected":N, "n\_ok":K, ...\}}) so a
        truncated cell is distinguishable from a complete one by inspection.
  \item \textbf{Cell-end barrier.} When the cell duration elapses the producer
        waits for outstanding helper acks before declaring the cell done and
        before any metric is computed. The shipped barrier uses a progressive
        idle timeout: exit when all expected acks arrive, OR no new ack appears
        for \SI{15}{\second} (idle quit), OR a \SI{180}{\second} hard cap fires.
        This guarantees the trajectory is complete (or its gaps are explicitly
        recorded) before downstream analysis.
        % HONESTY: the original barrier was a flat 30 s wall-clock wait; it
        % tail-missed 2/90 short ransom cells (the "trailing-helper barrier
        % race", Bug O). The progressive idle/cap barrier above is the fix that
        % brought the retry to 90/90 C6 pass.
  \item \textbf{Completeness validator (Δ-4).} A post-cell claim (C6) opens the
        trajectory, requires the final sentinel, and fails the cell if
        $n_{ok}/n_{pairs}$ falls below \SI{95}{\percent}; cells failing C6 are
        excluded from downstream aggregation. v1-era cells with no trajectory
        are reported ``not applicable'' (backward compatible).
  \item TODO: figure -- producer timeline (suspend / \texttt{pmemsave} / resume)
        overlapped with the async diff+delete helper and the cell-end barrier.
  \item TODO: figure -- disk-state diagram (prev/curr/in-flight) for
        delete-as-you-go vs the keep\_dumps end-of-cell-sweep reality.
\end{itemize}

%------------------------------------------------------------------
\section{Contributions in Detail}
\label{sec:contrib}
\begin{itemize}
  \item \textbf{C1 -- off-critical-path capture.} The only pausing work is the
        snapshot primitive; the page-change helper is fully asynchronous and
        joined by a cell-end barrier. Net: on-path cost reduced
        $\approx$\SI{2.5}{\second}~$\rightarrow$~$\approx$\SI{1.5}{\second}
        (proposal figure; verify vs Paper E).
  \item \textbf{C2 -- bounded transient disk (by design) + provable
        completeness.} Delete-as-you-go targets a $\approx$\SI{2}{}--\SI{3}{\gibi\byte}
        cap; ack files + structured exit codes + a seq/sentinel trajectory +
        the C6 validator make completeness an explicit, checkable property.
        % HONESTY: the disk cap is a target (unexercised in v3); the
        % completeness machinery IS exercised (C6 ran on all 90 cells).
  \item \textbf{C3 -- cost characterization + negative result.}
        \texttt{pmemsave} is the binding cost (Sec.~\ref{sec:eval}); concurrent
        multi-VM \texttt{pmemsave} is rejected with rationale
        (Sec.~\ref{sec:eval}, multi-VM).
\end{itemize}

%------------------------------------------------------------------
\section{Evaluation}
\label{sec:eval}
% Proposed results (bullets with REAL numbers from the v3 / timing runs):
\begin{itemize}
  \item \textbf{Snapshot cost dominates, and it is variable.} Over the v3 run
        (128 \texttt{ok} cells), the host snapshot cycle has median
        \SI{9.75}{\second} (range \SI{2.02}{}--\SI{23.57}{\second}).
        \texttt{pmemsave} is the dominant term: min \SI{0.76}{\second}, median
        \SI{7.06}{\second}, max \SI{21.57}{\second}, a median \SI{81}{\percent}
        of the cycle. Suspend ($\approx$\SI{0.35}{\second}) and resume
        ($\approx$\SI{0.05}{\second}) are negligible.
        % HONESTY: these are the campaign-level cost numbers (verified block).
        % Note Paper E's timing-instrumentation regime measured pmemsave at
        % ~0.77 s on a clean disk; the median 7.06 s here reflects the
        % production workload + disk-pressure conditions, not a clean-disk
        % floor. State both and explain the gap.
  \item \textbf{Async helper effect.} Moving the helper off-path removes the
        synchronous $\approx$\SI{1}{\second}/snapshot penalty of the B+3
        baseline; the proposal reports on-path
        $\approx$\SI{2.5}{\second}~$\rightarrow$~$\approx$\SI{1.5}{\second}.
        In the v3 run there were zero helper crashes (all
        \texttt{apf\_helper.log} files \SI{0}{\byte}) and zero non-zero exit
        codes.
        % HONESTY: per-pair helper compute (~1 s) is not separately logged in
        % v3; the 2.5->1.5 s figure is the design proposal's, not an isolated
        % v3 measurement. Measure helper compute directly before publication.
  \item \textbf{Trajectory completeness.} The C6 completeness gate passed on
        all production cells after the barrier fix: the initial run had 2/90
        cells fall below \SI{95}{\percent} (short ransom cells where the last
        pair's helper missed the flat \SI{30}{\second} barrier -- the
        helper did not crash; the sentinel recorded the gap honestly), and the
        progressive idle/cap barrier brought the retry to 90/90.
  \item \textbf{Trajectory length (points per cell).} min 4, median 30, max 99.
        Short cells are the stress case for the barrier (one missing pair is a
        large fraction of expected pairs).
  \item \textbf{Storage outcome (campaign-level).} The persisted APF trajectory
        is $\approx$\SI{4.7}{\mebi\byte} for a 90-cell campaign. Retaining all
        dumps for the full 2-D per-page pipeline projects to
        $\approx$\SI{5}{\tebi\byte} -- an estimated ratio of order $10^{6}$.
        % HONESTY: the ~5 TiB numerator is a PROJECTION for a campaign that was
        % never fully retained; the 4.7 MiB figure is for 90 cells. Present the
        % ratio as an estimate, not a measured ~10^6x superlative. This mirrors
        % Paper A's storage claim; keep the two papers consistent.
  \item \textbf{Transient disk: target vs reality.} The delete-as-you-go cap is
        $\approx$\SI{2}{}--\SI{3}{\gibi\byte}. In v3 (keep\_dumps on) deletion
        was deferred to a single end-of-cell sweep (cell notes record, e.g.,
        ``removed 29 dumps''), so transient disk scaled with trajectory length
        -- up to one \SI{1}{\gibi\byte} dump per snapshot, tens of \si{\gibi\byte}
        for a \SI{600}{\second} cell -- and the host-level free-disk swing for
        the run was \SI{59.5}{\gibi\byte}.
        % HONESTY (R-3, the paper's biggest open item): the 2-3 GiB cap was
        % NEVER exercised. Validating it requires a delete-as-you-go re-run;
        % that is the job of the Plan 05 snapshot-throughput proposal.
  \item \textbf{Negative result: concurrent multi-VM \texttt{pmemsave} (rejected).}
        Running multiple VMs and snapshotting them concurrently was evaluated
        against speed-up, honesty (cell independence), and production-fidelity.
        Rejected because:
    \begin{itemize}
      \item the host SSD is the bottleneck -- two simultaneous
            \SI{1}{\gibi\byte} \texttt{pmemsave} writes contend for disk
            bandwidth and serialize, so snapshot throughput does not improve and
            \texttt{pmemsave} timing becomes non-stationary (drifts);
      \item realistic wall-clock speed-up was only $\approx$\SI{30}{}--\SI{50}{\percent}
            (not the naive $3\times$), because the work is disk-bound not
            CPU-bound, and ``sampling-interleave'' saves nothing (same total
            number of writes);
      \item it breaks cell independence: overlapping cells share a
            ``host-contended'' signal, biasing any cross-cell comparison
            (correlated residuals), which is fatal for the per-cell measurements
            this engine exists to produce;
      \item it does not match the single-VM production/deployment model.
    \end{itemize}
        % Source: multi_vm_proposal.html (team verdict: DO NOT multi-VM for v2;
        % conditional-yes only for a different question -- multi-tenant claims,
        % 1000+ cell sweeps, or same-workload replicate parallelism with strict
        % CPU pinning + non-overlapping pmemsave windows).
  \item TODO: table -- snapshot-cycle decomposition (suspend / \texttt{pmemsave}
        / resume; min / median / max; pause fraction).
  \item TODO: table -- B+3 (sync) vs B+3.1 (async) on-path cost, disk peak, and
        completeness guarantee.
  \item TODO: figure -- transient-disk over a cell: delete-as-you-go target
        (flat $\approx$\SI{2}{\gibi\byte}) vs measured keep\_dumps growth.
\end{itemize}

%------------------------------------------------------------------
\section{Limitations}
\begin{itemize}
  \item \textbf{Delete-as-you-go was not exercised (R-3, primary honesty
        caveat).} The v3 campaign ran with \texttt{keep\_dumps} enabled on
        125/128 ok cells, which disables the helper's delete step. The
        $\approx$\SI{2}{}--\SI{3}{\gibi\byte} transient-disk cap is therefore a
        \emph{design target}, not a measured result; observed transient disk
        scaled with trajectory length and the host-level free swing reached
        \SI{59.5}{\gibi\byte}. Validating the cap (and shrinking
        \texttt{pmemsave}) is deferred to the Plan~05 snapshot-throughput
        proposal.
        % HONESTY: this is the paper's most important limitation -- keep it
        % first and explicit. Do not let the abstract/intro imply the cap was
        % measured.
  \item \textbf{\texttt{pmemsave} pauses the guest.} The snapshot primitive
        suspends the domain for the duration of the dump (median pause fraction
        high), so the engine is an offline/research capture tool, not a
        production-latency introspector. Live dirty-bit tracking is the natural
        optimization (would also remove the per-snapshot \SI{1}{\gibi\byte}
        write).
  \item \textbf{Per-cell vs host-level disk accounting.} Even with
        delete-as-you-go, the v3 evidence shows deletion can lag under
        \texttt{ionice} and block-level sweeps, so a per-cell quota (not just a
        per-cell steady-state argument) is needed before scaling to large or
        multi-host runs.
  \item \textbf{Barrier under short cells.} The completeness barrier is the
        delicate part for short trajectories (median 30 points, min 4): one
        late helper is a large fraction of expected pairs. The progressive
        idle/cap barrier fixed the observed cases but a hard cap remains a
        liveness fallback, not a correctness guarantee.
  \item \textbf{Single-host, single-SSD characterization.} The cost numbers and
        the multi-VM rejection are specific to one host with a single shared
        SSD; a different storage topology (e.g., per-VM NVMe) could change the
        multi-VM verdict (noted as a conditional re-evaluation).
  \item \textbf{Helper observability gap.} On success the helper log is empty
        (correct), but a crash signals only via a missing ack file; richer
        helper telemetry is future work.
\end{itemize}

%------------------------------------------------------------------
\section{Conclusion}
\begin{itemize}
  \item B+3.1 reduces whole-VM capture to a thin on-critical-path snapshot plus
        an asynchronous, delete-as-you-go signal helper, with a cell-end barrier
        that makes trajectory completeness an explicit, validated property.
  \item The binding cost is the hypervisor snapshot (\texttt{pmemsave}); the
        engine's central future lever is shrinking that primitive and finally
        exercising the delete-as-you-go disk bound -- both the subject of the
        Plan~05 snapshot-throughput proposal.
  \item Concurrent multi-VM \texttt{pmemsave} is rejected: on a single shared
        SSD it buys no throughput and costs cell independence. The
        signal-and-cost-model view of the same system is given in Paper A;
        the timing-instrumentation and interval-inertness evidence is in
        Paper E.
\end{itemize}

% ---------------------------------------------------------------------------
% Citation keys used above and their status in ../refs.bib:
%   garfinkel2003vmi            -- present (CANONICAL)
%   TODO_libvirt_pmemsave       -- present (placeholder; needs real pmemsave/virsh ref)
%   TODO_working_set_estimation -- present (placeholder)
%   TODO_ksm_dedup              -- present (placeholder)
% Still TO ADD to ../refs.bib before submission:
%   - a snapshot/checkpoint cost-study citation (for the pmemsave-overhead framing)
%   - an OS page-cache / write-back behavior citation (for the dirty-page coupling)
%   - (optional) a prior streaming / online-diff capture system, to position
%     delete-as-you-go.
% ---------------------------------------------------------------------------

\bibliographystyle{IEEEtran}
\bibliography{../refs}
\end{document}
